% ============================================================================
% Abstract — rewritten 2026-10-19 on the final numbers (baseline re-run under
% the comparison protocol, one scoring rule, TruthfulQA judge v2).
% Every number traces to outputs/analysis/chm_judged.json,
% outputs/analysis/paired_baseline_tests_final.json and the chapter-five tables
% built by scripts/make_final_results.sh.
% ============================================================================

\section*{Abstract}

Large language models give confident but false answers, and a user cannot tell from the answer which ones to trust: asked for supporting literature, chatbots have returned references that do not exist \cite{alkaissi2023artificial}, for up to $91.4$ percent of the references they generate in systematic-review searches \cite{info:doi/10.2196/53164}. This thesis proposes \emph{Confidence-Aware Episodic Memory with Self-Improvement (CAEM)}, a system built around a small open generator. A router sends each question to a verified memory, to zero-shot generation or to retrieval-augmented generation, with a classifier that decides when retrieval is likely to help. A verifier combines nine signals of internal confidence, agreement across sampled answers, support from retrieved passages and question match into a calibrated probability of correctness, and a four-outcome decision stores, defers, abstains on or discards each answer. Every answer reaches the user with that probability and a status tag (verified, provisional, conflicting or insufficient evidence). Verified answers are kept in memory, re-checked at every cycle, and used to fine-tune a low-rank adapter, with a retention guard that rolls the adapter back when held-out accuracy falls. Two results are proven: storage makes the stored pool purer than the generator's raw output exactly when the verifier's balanced accuracy exceeds one half, and the pool's precision is bounded below by the calibration precision when new questions resemble the calibration data.

CAEM was run for six cycles with Qwen-2.5-3B-Instruct on five benchmarks (FEVER, TriviaQA and CommonsenseQA for training; TruthfulQA and StrategyQA held out) and compared with eight baselines under one protocol: the same questions, the same prompt, one scoring rule and the same hallucination instrument. The architecture alone added almost no accuracy (pooled exact match $0.517$ against $0.511$ for the zero-shot baseline); self-improvement raised it to $0.535$ after two cycles, after which it eroded to $0.508$, and no difference from the zero-shot baseline is significant on any benchmark. CAEM is level with strong prompting baselines, significantly better than the retrieval baselines on several benchmarks, and level on accuracy with a fine-tuned baseline that used about eleven times as many gold labels, which reached a lower hallucination metric. Its measured contribution is selectivity: the status tag predicts correctness, with verified answers right $65$ to $75$ percent of the time against $51$ percent for the zero-shot baseline. The gains are concentrated where the verifier can check answers against retrieved passages, and after self-improvement the verifier scored more of the wrong answers highly (confident confabulation $0.305$ against $0.208$). Of four pre-specified hypotheses, two are supported, one partially and one not.

\vspace{0.5cm}
\noindent\textbf{Keywords:} hallucination; episodic memory; calibrated confidence; selective prediction; verifier; self-improvement; retrieval-augmented generation; low-rank adaptation; large language models.
\pagebreak
